\subsection{Inverting Amplifier}\label{sec:amp_inverting}

\subsubsection{Overview}
An op-amp with \(V_+\) at ground, the input through \(G_1\) into \(V_-\) and
feedback through \(G_2\). Every functional block uses the op-amp model of
Section~\ref{sec:devices_opamp}: \(V_{out} = A(V_+ - V_-)\), with \(A = A_0\) in Q5 and
\(A = A(s) = \frac{A_0 p_1}{s + p_1}\) in Q4; the ideal result is the limit \(A \to \infty\).

\subsubsection{Schematic}
\begin{circuitfig}
  \draw (0,0) node[op amp](A){};
  \draw (A.-) -- ++(-0.6,0) coordinate(n) node[circ]{} node[below left]{$V_-$} to[european resistor, l_={$G_1$}] ++(-2.4,0) node[circ]{} node[left]{$V_{in}$};
  \draw (n) |- (-0.5,1.6) to[european resistor, l={$G_2$}] (1.6,1.6) -| (A.out);
  \draw (A.+) -- ++(-0.6,0) node[ground]{};
  \draw (A.out) to[short, -o] ++(0.8,0) node[right]{$V_{out}$};
\end{circuitfig}

\subsubsection{Transfer Function}
\begin{enumerate}
  \item \textbf{Large Signal, Time Domain, DC Input Analysis (Q5)}: \(A = A_0\),
    \(V_+ = 0\), no current into the op-amp inputs.
    \begin{align*}
      V_{OUT} &= A_0(V_+ - V_-) = -A_0 V_- \hspace{1cm} \Rightarrow \hspace{1cm} V_- = -\frac{V_{OUT}}{A_0} \tag{1}
    \end{align*}
    Applying KCL at \(V_-\):
    \begin{align*}
      (V_- - V_{IN}) G_1 + (V_- - V_{OUT}) G_2 &= 0 \\
      V_-(G_1 + G_2) - V_{IN} G_1 - V_{OUT} G_2 &= 0 \\
      -\frac{V_{OUT}}{A_0}(G_1 + G_2) - V_{OUT} G_2 &= V_{IN} G_1 \hspace{1cm} [\text{from (1)}] \\
      -V_{OUT}\left(G_2 + \frac{G_1 + G_2}{A_0}\right) &= V_{IN} G_1 \\
      \frac{V_{OUT}}{V_{IN}} &= -\frac{G_1}{G_2 + \frac{G_1 + G_2}{A_0}} \tag{2}
    \end{align*}
    valid while \(V_{OL} \leq V_{OUT} \leq V_{OH}\) (outside it the output clips).
  \item \textbf{Small Signal, Frequency Domain, AC Input Analysis (Q4)}: the same KCL with
    \(A = A(s)\), so (2) becomes
    \begin{align*}
      H(s) = \frac{V_{out}}{V_{in}} &= -\frac{G_1}{G_2 + \frac{(G_1 + G_2)(s + p_1)}{A_0 p_1}} \hspace{1cm} \left[\tfrac{1}{A(s)} = \tfrac{s + p_1}{A_0 p_1}\right] \\
      &= -\frac{G_1 A_0 p_1}{G_2 A_0 p_1 + (G_1 + G_2)(s + p_1)} \\
      &= -\frac{G_1 A_0 p_1}{(G_1 + G_2)s + p_1(G_1 + G_2 + A_0 G_2)} \\
      &= -\frac{G_1}{G_1 + G_2}\cdot\frac{A_0 p_1}{s + p_1\left(1 + A_0\frac{G_2}{G_1 + G_2}\right)}
    \end{align*}
    With the feedback factor \(F = \frac{G_2}{G_1 + G_2}\) (Analysis, 5):
    \begin{align*}
      H(s) &= -\frac{G_1}{G_1 + G_2}\cdot\frac{A_0 p_1}{s + p_1(1 + A_0 F)} \tag{3}
    \end{align*}
\end{enumerate}

\subsubsection{Analysis}
\begin{enumerate}
  \item \textbf{Poles (\(D(s) = 0\))}:
    \begin{align*}
      s_p + p_1(1 + A_0 F) &= 0 \hspace{1cm} \Rightarrow \hspace{1cm} s_p = -p_1(1 + A_0 F) \\
      f_p &= \frac{|s_p|}{2\pi} = \frac{p_1(1 + A_0 F)}{2\pi} \approx F \cdot \frac{A_0 p_1}{2\pi} = F \cdot GB \hspace{1cm} [A_0 F \gg 1]
    \end{align*}
    Feedback moves the op-amp pole up by \(1 + A_0 F\).
  \item \textbf{Zeroes (\(N(s) = 0\))}: \(N(s) = -\frac{G_1}{G_1 + G_2} A_0 p_1 \neq 0\): no finite zero.
  \item \textbf{DC Gain (\(H(0)\))}: \(s = 0\) in (3),
    \begin{align*}
      H(0) &= -\frac{G_1}{G_1 + G_2}\cdot\frac{A_0}{1 + A_0 F}
            = -\frac{G_1}{G_2}\cdot\frac{G_2}{G_1 + G_2}\cdot\frac{A_0}{1 + A_0 F} = -\frac{G_1}{G_2}\cdot\frac{A_0 F}{1 + A_0 F} \\
      H(0) &\xrightarrow{A_0 \to \infty} -\frac{G_1}{G_2} = -\frac{R_2}{R_1}
    \end{align*}
    The gain error is \(1 - \frac{A_0 F}{1 + A_0 F} = \frac{1}{1 + A_0 F}\).
  \item \textbf{Gain-Bandwidth Product (\(GBW\))}:
    \begin{align*}
      GBW &= |H(0)| \cdot f_p = \frac{G_1}{G_1 + G_2}\cdot\frac{A_0}{1 + A_0 F}\cdot\frac{p_1(1 + A_0 F)}{2\pi}
           = \frac{G_1}{G_1 + G_2}\cdot\frac{A_0 p_1}{2\pi} = (1 - F)\, GB
    \end{align*}
    (\(\frac{G_1}{G_1 + G_2} = 1 - F\)): less than \(GB\), because the signal sees \(\frac{R_2}{R_1}\) while
    the loop sees \(\frac{1}{F} = 1 + \frac{R_2}{R_1}\).
  \item \textbf{Feedback Factor (\(F\))}: \(V_{in} = 0\), KCL at \(V_-\):
    \begin{align*}
      V_- G_1 + (V_- - V_{out}) G_2 &= 0 \hspace{1cm} \Rightarrow \hspace{1cm} V_- = \frac{G_2}{G_1 + G_2} V_{out} \\
      F &= \frac{G_2}{G_1 + G_2} = \frac{\frac{1}{R_2}}{\frac{R_1 + R_2}{R_1 R_2}} = \frac{R_1}{R_1 + R_2}
    \end{align*}
  \item \textbf{Impedances}:
    \paragraph{At Input:} \(V_- = -\frac{V_{out}}{A} \to 0\) (virtual ground), so the source sees \(G_1\) to ground:
    \begin{align*}
      I_{in} &= (V_{in} - 0) G_1 \hspace{1cm} \Rightarrow \hspace{1cm} Z_{in} = \frac{V_{in}}{I_{in}} = \frac{1}{G_1} = R_1
    \end{align*}
    \paragraph{At Output:} \(V_{in} = 0\), test voltage \(V_{test}\) at the output, op-amp output
    resistance \(R_o\) in series with its source \(-A V_-\), \(V_- = F V_{test}\) (\(R_1 + R_2 \gg R_o\)):
    \begin{align*}
      I_{test} &= \frac{V_{test} - (-A F V_{test})}{R_o} = \frac{V_{test}(1 + AF)}{R_o} \hspace{1cm}
      Z_{out} = \frac{V_{test}}{I_{test}} = \frac{R_o}{1 + AF}
    \end{align*}
  \item \textbf{Output Swing (Max)}: \(V_{OL} \leq -\frac{R_2}{R_1} V_{IN} \leq V_{OH}\), so
    \(-\frac{R_1}{R_2} V_{OH} \leq V_{IN} \leq -\frac{R_1}{R_2} V_{OL}\).
  \item \textbf{Slew Rate (Large Signal, Time Domain, AC Input Analysis (Q6))}: a sine output
    \(V_p \sin(2\pi f t)\) has the largest slope \(2\pi f V_p\), so it is undistorted while
    \(2\pi f V_p \leq SR\), i.e. \(f \leq \frac{SR}{2\pi V_p}\).
  \item \textbf{Offset and Bias Currents (Large Signal, Time Domain, DC Input Analysis (Q5))}: \(V_{IN} = 0\),
    offset \(V_{OS}\) in series with \(V_+\), bias currents \(I_{B+}\), \(I_{B-}\) into the inputs, a resistor \(R_p\) from
    \(V_+\) to ground. Ideal op-amp, \(V_- = V_+ = V_{OS} - I_{B+} R_p\); applying KCL at \(V_-\):
    \begin{align*}
      \frac{V_-}{R_1} + \frac{V_- - V_{OUT}}{R_2} + I_{B-} &= 0 \\
      V_{OUT} &= V_-\left(1 + \frac{R_2}{R_1}\right) + I_{B-} R_2 = \left(1 + \frac{R_2}{R_1}\right)(V_{OS} - I_{B+} R_p) + I_{B-} R_2
    \end{align*}
    \(V_{OS}\) is amplified by the noise gain \(\frac{1}{F} = 1 + \frac{R_2}{R_1}\), not by \(\frac{R_2}{R_1}\). With
    \(I_{B+} = I_{B-} = I_B\) and \(R_p = R_1 \parallel R_2 = \frac{R_1 R_2}{R_1 + R_2}\):
    \begin{align*}
      \left(1 + \frac{R_2}{R_1}\right) I_B \frac{R_1 R_2}{R_1 + R_2} &= \frac{R_1 + R_2}{R_1}\cdot\frac{R_1 R_2}{R_1 + R_2} I_B = I_B R_2 \hspace{1cm} \Rightarrow \hspace{1cm} V_{OUT} = \frac{V_{OS}}{F} \pm I_{OS} R_2
    \end{align*}
    the bias currents cancel, leaving only the offset current \(I_{OS} = |I_{B+} - I_{B-}|\).
  \item \textbf{Noise}: an input referred op-amp noise density \(\overline{V_n^2}\) appears at the output times
    \(\frac{1}{F^2}\); each resistor's \(4kTR\) adds: \(\overline{V_{o}^2} = \frac{\overline{V_n^2}}{F^2} + 4kTR_1\left(\frac{R_2}{R_1}\right)^2 + 4kTR_2\) (per Hz,
    Section~\ref{sec:theory_general_noise}).
  \item \textbf{Advantages \& Limitations}: gain set by a resistor ratio, input
    common mode fixed at 0 (no common mode error); but \(Z_{in} = R_1\) loads the source,
    the output is inverted, and \(GBW = (1 - F) GB\).
\end{enumerate}

\subsubsection{Design}
\begin{enumerate}
  \item Choose \(R_1 = Z_{in}\) required by the source.
  \item \(R_2 = |H(0)| R_1\) for the required gain.
  \item Check \(f_{-3dB} = F \cdot GB = \frac{R_1}{R_1 + R_2} GB\) meets the bandwidth, and the
    gain error \(\frac{1}{1 + A_0 F}\) meets the accuracy.
\end{enumerate}

\subsubsection{Question}
\(R_1 = 1\,k\Omega\), \(R_2 = 10\,k\Omega\), \(A_0 = 10^5\), \(GB = 1\,\)MHz. Find the exact DC gain,
the bandwidth and \(GBW\).
\paragraph{Solution:}
\begin{align*}
  F &= \frac{R_1}{R_1 + R_2} = \frac{1\,k}{1\,k + 10\,k} = \frac{1}{11} \hspace{1cm} A_0 F = \frac{10^5}{11} = 9090.9 \\
  H(0) &= -\frac{R_2}{R_1}\cdot\frac{A_0 F}{1 + A_0 F} = -10 \cdot \frac{9090.9}{9091.9} = -10 \cdot 0.99989 = -9.9989 \\
  f_p &\approx F \cdot GB = \frac{1\,MHz}{11} = 90.9\,kHz \\
  GBW &= (1 - F)\, GB = \frac{10}{11} \cdot 1\,MHz = 909\,kHz \hspace{1cm} [= 9.9989 \cdot 90.9\,kHz]
\end{align*}
